\section{Morpheus} \label{sec:morpheus}
\begin{figure}
  \centering
  \makebox[\textwidth]{
    \includegraphics[width=1.3\textwidth]{figures/03_morpheusEvasion/morpheusFrames/crop/merge.png}
  }
  \caption{
    The motion is generated over $T_f=\SI{1.5}{\second}$ with $N=48$ shooting nodes for a head reduction of $\Delta h=\SI{0.5}{\metre}$.
    We visualize the every eighth state.
    As demanded by the stability objective function the red marked \acrshort{zmp} inside the \acrshort{pos} remains mostly centered.
  }
  \label{fig:morpheus:results:banner}
\end{figure}
We start the development of our collision avoidance motion generator by focusing on objects, which are targeted towards the head.
An intuitive and easy solution is to move the head downwards and out of the way so that the object passes right over.
\\
The working title of this first generator is \textit{Morpheus} after the famous character from "The Matrix".
Like the character in the movie, this formulation is the steppingstone for our more advanced \textit{Neo} formulation in \cref{sec:neo}.
\par
\textit{Morpheus} idea is to formulate a well rounded whole-body locomotion optimal control problem, that is able to reliably generate motions for the humanoid H1-2 in the context of object avoidance.
We start of by discussing the trade-offs taken during model selection (\cref{sec:morpheus:model}) and than move towards describing the exact problem formulation in \cref{sec:morpheus:formulation}.
There we also present the choice of states, controls, dynamics, constraints, and objective function.
An investigation of the available hardware limits of the H1-2 rounds up the formulation.
\\
At last we discuss the shortcomings and achievements of our \textit{Morpheus} implementation on the example of a motion, that was generated for a head height reduction of $\Delta h = \SI{0.5}{\metre}$.
Additionally, we assess the problems posed by the robots hardware limitations and compare them to the capabilities of a human body.

\subsection{Robot Model} \label{sec:morpheus:model}
As every model is an approximation of the real world, we have to model it in a way to capture the essential details for our problem.
A relatively simple decision to take is the use of a floating base model instead of a mounted model.
The floating base joint vector
\begin{equation}
  \bm q =
  \begin{bmatrix}
    \bm p_\text{world} \\
    \bm r_\text{world} \\
    \bm q_\text{joints}
  \end{bmatrix}
  \label{eq:morpheus:joints}
\end{equation}
holds not only the actual joints $\bm q_\text{joints} \in \mathbb{R}^{n_j}$, but also the world position $\bm p_\text{world} \in \mathbb{R}^3$ and orientation $\bm r_\text{world} \in \mathbb{R}^3$.
In some applications the orientation is defined as quaternion to provide singularity free transitions.
This however would lead to complications in calculating the joint velocities as not every rigid-body library supports quaternion velocities.
Therefore, we decided to stick to euler angles for an easy to use system.
In total our joint vector has the size of $n_q=6+n_j$.
\\
The goal of our motion generator is to reduce the head height while remaining in bipedal stance.
To achieve this, we do not need to employ joints outside the sagittal plane and can reduce the problem size therefore to seven joints: two ankle, knee, and hip joints plus the torso joint.

\begin{figure}
  \centering
  \includegraphics[height=8cm]{figures/03_morpheusEvasion/lumped_model.png}
  \caption{The movable segments of the defined robot model. Excluding the torso it is restricted to sagittal motions.}
  \label{fig:morpheus:lumped}
\end{figure}
The resulting movable segments are marked in \cref{fig:morpheus:lumped}:
Purple marks the torso, blue the thigh, green the shank and orange the feet.
An H1-2 model, which includes all joints is provided by Unitree.
With the use of our rigid-body dynamics library of choice, we are able to fixate all not required joints, so that they can be considered as solid links in the following.
After we finished selecting our robot model, we move to defining the optimal control problem.

\subsection{Optimal Control Formulation} \label{sec:morpheus:formulation}
To achieve highly dynamic motion synthesis, we select for our controls $\bm u$ the joint torques $\bm \tau$ of the robot.
These controls are a common choice and were for example used to generate somersault motions, like Mombaur et al. \cite{mombaur_2005} showed.
Additionally, it gives the open-loop controller as much controller as possible and closely ties to the real robots control inputs.
As a result the states $\bm x$ of our system are composed from joint positions $\bm q$ and velocities $\bm{\dot q}$, which boils done to:
\begin{subequations}
  \label{eq:morpheus:dyn}
  \begin{align}
    \bm x &=
    \begin{bmatrix}
      \bm{q} \\ \bm{\dot q}
    \end{bmatrix} \label{eq:morpheus:dyn:x} \\
    \bm u &=
    \bm{\tau} \label{eq:morpheus:dyn:u}
  \end{align}
\end{subequations}
With these we are able to formulate the objective function, explicit system dynamics, as well as constraints.
The following presents the full formulation of our optimal control problem:
\begin{subequations}
  \label{eq:morpheus:constraints}
  \begin{align}
    & \qquad \textbf{Objective Function} \nonumber \\
    \min_{\bm x(\cdot),\bm u(\cdot)} \quad
      & \int_0^{T_f} \frac 1 2 \left(\omega_s l_s(\bm x)^2 + \omega_r l_r(\bm{u})^2 \right) dt
      + \frac 1 2 \omega_h m_h(\bm x_f)^2 \label{morpheus:constraints:obj} \\
    % Dynamics
    & \qquad \textbf{System Dynamics} \nonumber \\
    \textrm{s.t.} \quad & \bm{\dot x} =
    \begin{bmatrix}
      \bm{\dot q} \\
      \text{cons. FD}(\bm q,\bm{\dot q},\bm{\tau})
    \end{bmatrix} \label{eq:morpheus:constraints:f} \\
    % Constraints
    & \qquad \textbf{Initial Pose Constraint} \nonumber \\
    & \begin{bmatrix} \bm q_0 \\ 0 \end{bmatrix} = \bm x(0) & & \label{eq:morpheus:constraints:initialPose} \\
    & \qquad \textbf{Dodge Constraint} \nonumber \\
    & \bm p_\text{head,f}|_z \leq \bm p_\text{head,0}|_z - \Delta h & & \label{eq:morpheus:constraints:dodge} \\
    & \qquad \textbf{Stability Constraint} \nonumber \\
    & \text{ZMP} \in \gamma \text{PoS} & & \forall t \in [0,T_f) \label{eq:morpheus:constraints:stability} \\
    & \qquad \textbf{Joint Constraints} \nonumber \\
    & \bm{\underline q} \leq \bm{q} \leq \bm{\overline q} & & \forall t \in [0,T_f] \label{eq:morpheus:constraints:joint}\\
    & -\bm{\tau}_\text{max} \leq \bm u \leq \bm{\tau}_\text{max} & & \forall t \in [0,T_f)\label{eq:morpheus:constraints:torque}\\
    & \qquad \textbf{Friction Constraints} \nonumber \\
    & 0 \leq (\mu F_j|_z)^2-\|\bm{F}_j\|^2_2 & \forall j \in \{ \text{left}, \text{right} \} \qquad & \forall t \in [0,T_f) \label{eq:morpheus:constraints:F} \\
    & 0 \leq F_j|_z & \forall j \in \{ \text{left}, \text{right} \} \qquad & \forall t \in [0,T_f) \label{eq:morpheus:constraints:Fz} \\
    & \qquad \textbf{Terminal Pose Constraint} \nonumber \\
    & -\bm v_\text{stat.} \leq \bm{\dot q_f} \leq \bm v_\text{stat.} \label{eq:morpheus:constraints:terminal}
  \end{align}
\end{subequations}
During the next sections we will walk through the components of $\cref{eq:morpheus:constraints}$ step by step.
However before we pivot into the details of theoretical formulations, we need to discuss some key points of the implementation.
\par
Joint accelerations can be calculated from positions, velocities, and torques using constraint forward dynamics.
We settled for a solution of the in \cref{eq:consDyn:default} defined system, as it is already implemented in the rigid-body dynamics library Pinocchio \cite{pinocchioweb,carpentier2019pinocchio}.
\\
To fully utilize the strength of our selected optimal control library Acados \cite{verschueren_2021} we employ automatic or also known as algorithmic differentiation in the form of CasADi \cite{andersson2019} expressions.
Pinocchio \cite{carpentier2019pinocchio,pinocchioweb} as well as \acrfull{rbdl} \cite{Felis2016} are able to produce CasADi expressions for a given \verb|urdf|-model.
Both are based on Featherstones methods and equations \cite{featherstone_2008}.
We chose Pinocchio as it is the faster \cite{carpentier2019pinocchio} and more commonly used library.
Another advantage is the availability of visualizers and spread of available algorithms it overs.
As a result Pinocchio is used for all our dynamics and kinematics equations of the H1-2, as well as the contact formulation mentioned before.
\\
Because Pinocchio has specific restrictions on how to define contact points, while our robot model is moving solely in the sagital plane, we land on a singular \acrshort{kkt} matrix $K(\bm q)$ in \cref{eq:consDyn:default:KKT}.
Carpentier et al. \cite{carpentier_2021} present a reformulation with the proximal operator as a solution to this problem, like shown in \cref{eq:consDyn:prox:operator,eq:consDyn:prox:algo}.
This makes the \acrshort{kkt} matrix invertable, but introduces an inaccuracy (\cref{eq:consDyn:prox:KKT}).
As the automatic differentiation requires a closed form of our contact formulation, we unroll a specified amount of iterations for \cref{eq:consDyn:prox:algo}, which will result in a suitable approximation.
\\
After clarifying the system dynamics and framework choices we move towards the composited objective function and limiting constraints in \cref{eq:morpheus:constraints}.

\subsubsection{Objective Functions}
The objective function represents the features we would like to achieve, when formulating an optimal control problem.
For our humanoid collision avoidance robot the main feature is to perform a matrix style bullet dodge, where the robot leans back, while lowering its head height.
In order to avoid the head targeted approaching object we require the robot to lower its height more than $\Delta h$.
As this is a hard limit, it is formulated as a constraint in \cref{eq:morpheus:constraints:dodge}.
\par
To resemble a full matrix style dodge, we would like to maximize the lean back of H1-2.
Therefore, the head position $\bm p_\text{head}$ needs to be pushed into the negative $x$ direction.
For acados to work properly we have to formulate our objective function in the following way:
\begin{equation}
  \min_{\bm x(\cdot),\bm u(\cdot)} \quad
  \int_0^{T_f} \frac 1 2 \sum^{n_l}_{i}{\left(\omega_i l_i(\bm x, \bm u)^2 \right)} dt
  + \frac 1 2 \sum^{n_m}_{j}{\left(\omega_j m_j(\bm x_f)^2\right)}
  \label{eq:morpheus:acados:obj}
\end{equation}
As the reader might have noticed this forces every Lagrange or Mayer term to be defined in the range of $[0,+\inf]$.
So therefore the intuitive solution of directly formulating the Mayer term 
\begin{equation}
  m_h(\bm x_f) = \bm{p_\text{head}}|_x
  \label{eq:morpheus:naive}
\end{equation}
results in a quadratic function with its optimum at $p_\text{head}|_x = 0$.
\Cref{fig:morpheus:objHead:naiv} displays $m_h(\bm x_f)^2$ and marks the values for when the head is over the front $\overline{PoS}|_x$ and back $\underline{PoS}|_x$ of the \acrshort{pos}.
\begin{figure}
  \centering
    \begin{subfigure}[b]{0.45\textwidth}
    \includegraphics[draft,width=\textwidth,height=5cm]{figures/03_morpheusEvasion/naive/naive.pdf}
    \caption{Naive Approach}
    \label{fig:morpheus:objHead:naiv}
  \end{subfigure}
  \hfill
  \begin{subfigure}[b]{0.45\textwidth}
    \includegraphics[width=\textwidth]{figures/03_morpheusEvasion/sigmoid/sigmoid.pdf}
    \caption{Logistic Function Approach}
    \label{fig:morpheus:objHead:sigmoid}
  \end{subfigure}
  \caption{
    The resulting value range are displayed for (a) the in \cref{eq:morpheus:sigmoid} and (b) \cref{eq:morpheus:sigmoid} after squaring them.
    It is easy to see, that for (a) the objective function would pull the robot into the center, while (b) emphasizes a lean over.
    For better clarity of direction the values of at the edge of the \acrshort{pos} are marked.
  }
  \label{fig:morpheus:objHead}
\end{figure}
\\
In order to project negative $x$-coordinates to near zero and positive values to one we employ a sigmoid curve in form of the logistic function:
\begin{equation}
  s(x,k_0,k_1,x_0) = \frac{k_0}{1+e^{-k_1(x-x_0)}}
  \label{eq:morpheus:sigmoid}
\end{equation}
It reassembles the characteristic s-shape of a sigmoid function, by defining the supremum $k_0 = 1$, steepness $k_1 = 30$ and cross point $x_0 = 0$.
The resulting function is shown together with markings for the \acrshort{pos} in \cref{fig:morpheus:objHead:sigmoid}.
Through the logistic function we are able to apply a consistent penalty for leaned forward poses.
A shift of the sigmoid curve towards the \acrshort{pos} center is not advised, as this would ignore the natural bias of the humanoid robot, as its \acrshort{com} in the default position is at $\textbf{CoM}_0|_x = 0$.
This is the cause of the robots feet being not symmetrical in $x$-direction, but resemble more the ones of a human, where the ankle joint is sitting towards the heel of the foot.
\\
By utilizing the logistic function $s(x)$ from \cref{eq:morpheus:sigmoid} we formulate the Mayer Objective Function $m_h$.
Together with the Largrange terms $l_s$ and $l_r$ all objectives can be written as follows:
\begin{subequations}
  \label{eq:morpheus:obj}
  \begin{align}
    m_h(\bm x_f) &= s(p_\text{head}|_x) \label{eq:morpheus:obj:head} \\
    l_s(\bm x) &= 
      \left(\frac{\textbf{ZMP}|_x - \textbf{PoS}_\text{center}|_x}{0.5\cdot\text{PoS}_\text{height}}\right)^2 
      + \left(\frac{\textbf{ZMP}|_y - \textbf{PoS}_\text{center}|_y}{0.5\cdot\text{PoS}_\text{width}}\right)^2
      \label{eq:morpheus:obj:stability} \\
    l_r(\bm u) &= \left\lVert \frac{\bm{\tau}}{\bm{\tau}_\text{max}} \right\rVert^2_2 \label{eq:morpheus:obj:regularization}
  \end{align}
\end{subequations}
The stability Lagrange objective function $l_s$ uses the \acrshort{zmp} from \cref{eq:dynStab:ZMP}.
Like we already mentioned in \cref{sec:fundamentals:dynamicStability} the \acrshort{zmp} has to be well inside the \acrshort{pos} to guarantee stability.
Generally the \acrshort{zmp} exists only inside the \acrshort{pos} and during none slipping contact.
The later is ensured by constraining the allowed friction \cref{eq:morpheus:constraints:F,eq:morpheus:constraints:Fz} with a friction coefficient of $\mu = 0.5$.
\Cref{eq:morpheus:constraints:Fz} further ensures that no "pulling" contact forces are exploited, as this is mathematically possible with the formulation in \cref{eq:consDyn:default}, but would not be feasible in terms of physics.
To assure the existence of the \acrshort{zmp} we clamp it to the \acrshort{pos}.
\Cref{eq:morpheus:constraints:stability} amplifies the restriction by reducing the \acrshort{pos} with the factor $\gamma$.
This is a common practice to account for uncertainties in the execution of the motion later in the pipeline.
\\
Additionally, we push the \acrshort{zmp} into the middle of the \acrshort{pos} by minimizing for the quadratic normalized distance between center and \acrshort{zmp}.
Because we are optimizing over the full time we model it as the Lagrange objective function $l_s$, which can be found in \cref{eq:morpheus:obj:stability}.
\\
The last part of our objective function is a regularization term $l_r$.
Like the other components it was normalized to offer a value range of $\left[0,1\right]$.
For the regularization the individual torque limits of the joints are used $\bm \tau_\text{max}$.
These are the same we will later use to constraint the available torque.
\par
As a result we are able to weigh the objective function terms in order of importance.
Stability and lean are the main optimization goals and should therefore receive the same weighting.
To balance the fact out that the Lagrange function is integrated over the full duration, we scale the lean weight of our Mayer term by the amount of shooting nodes $N$.
Regularization, as the name already suggests, should remain small and not hinder the achievement of the other goals, its weight $\omega_r$ remains low.
\\
The exact weights are as follows:
\begin{subequations}
  \begin{align}
    \omega_s &= 10^{-4} & \text{stability} \nonumber \\
    \omega_h &= N & \text{head / lean} \nonumber \\
    \omega_r &= 10^{-8} & \text{regularization} \nonumber
  \end{align}
  \label{eq:morpheus:weights}
\end{subequations}
\par
We already covered some elementary constraints of the optimal control problem.
The next section fills in the gaps by describing the rational behind the initial pose \cref{eq:morpheus:constraints:initialPose}, evasion constraint \cref{eq:morpheus:constraints:dodge}, hardware limits \cref{eq:morpheus:constraints:joint} and terminal constraint \cref{eq:morpheus:constraints:terminal}.

\subsubsection{Additional Constraints}
Additional to the already discussed stability and friction constraints (\cref{eq:morpheus:constraints:dodge,eq:morpheus:constraints:F,eq:morpheus:constraints:Fz}) there are initial as well as terminal pose constraints, plus hardware limits in form of joint angle and torque constraints.
\\
The initial pose constraint \cref{eq:morpheus:constraints:initialPose}, ensures, that the solver is generating a trajectory, where the robot starts from a stationary predefined pose.
The bespoken pose is defined with slight bend knees ($\SI{0.4}{\radian}$) to avoid the singularity prone straight knee pose.
It also enables the robot to dampen rapid changes in $z$-direction over the ankle, knee, and hip joint.
\par
For an effective usage of the robots capabilities it is important to define the available resources as constraints.
In our formulation these are set in \cref{eq:morpheus:constraints:joint}.
As Unitree provides an officially supported \verb|urdf| model, we use the joint angle limits from there.
\\
However, when taking a closer look at the torque, or as it is called in the \verb|urdf| file, \verb|effort| they used significantly reduced values, compared to their values on the official Website \cite{unitreeH1-2}.
The intention is likely to set a safety margin for controllers and especially for learning based methods.
As our goal is the synthesis of motions right on the edge of the hardware capabilities, we have to trust the executing \acrfull{wbc} to limit torques and protect them from overworking the motors.
Still, we would like the motion generator to use as much of the available torque as possible, driving the robot to its limits, will not overshooting.
The result are torque limits we directly collected from the Unitree website \cite{unitreeH1-2} with addition to some of our own analysis of the robots internals.
\\
\begin{table}[H]
  \centering
  \begin{tabularx}{0.3\textwidth}{rc}
    \toprule
    Joint Name & Limits [Nm] \\
    \midrule
    Torso & 220 \\
    Hip Pitch & 220 \\
    Knee & 360 \\
    Ankle Pitch & \textit{130} \\
    \bottomrule
  \end{tabularx}
  \caption{Torque Limits}
  \label{tab:morpheus:limits}
\end{table}
\Cref{tab:morpheus:limits} displays the resulting values.
While the torso, hip pitch and knee joints torque limits can be taken directly from the website, the two ankle joints, pitch and roll, are driven by two motors with a limit of $\SI{70}{\newton \metre}$ each.
Therefore, the joints torque limits should not be set independently, as this results in over or underestimating the available torques.
Instead, a combined solution needs to be introduced.
\\
We support this claim with images from the internal of the H1-2 in the BioRobotics Lab at KIT.
\begin{figure}
  \centering
  \begin{subfigure}[b]{0.34\textwidth}
    \centering
    \includegraphics[width=\textwidth]{figures/03_morpheusEvasion/ankle_top.jpeg}
    \caption{Top View}
    \label{fig:morpheus:ankle:top}
  \end{subfigure}
  \hfill
  \begin{subfigure}[b]{0.55\textwidth}
    \centering
    \includegraphics[width=\textwidth]{figures/03_morpheusEvasion/ankle_side.jpeg}
    \caption{Side View}
    \label{fig:morpheus:ankle:side}
  \end{subfigure}
  \caption{
    Those pictures were taken from the H1-2 in the BioRobotics Lab of the KIT.
    They show the internals of the pitch and roll joint in the ankle.
    The driving lever arms $a$ were measured with $\SI{30}{\milli \metre}$ (a), while the one at the joint $b$ is $\SI{26}{\milli \metre}$ resulting in an approximate conversion rate of $0.87$.
  }
  \label{fig:morpheus:ankle}
\end{figure}
\Cref{fig:morpheus:ankle} shows the two motors, which drive the pitch and roll ankle joints together.
Because we solely require torques in the pitch direction for sagittal plane motions, we can contribute the whole performance of the motors towards this joint.
The differing lever arms $a$ and $b$ at the motor and ankle lead to a reduced maximum torque of $\SI{130}{\newton \metre}$.
With information about the limitation of our robotics hardware we define the constraint \cref{eq:morpheus:constraints:joint}.
\par
For the terminal pose constraint, the robot should stop in a quasi-static pose.
As long as the solution fulfills the dodge constraint (\cref{eq:morpheus:constraints:dodge}), the final poses joint angels $\bm q$ are variable.
This means our terminal constraint restricts solely the joint velocities at the end time $T_f$.
\\
We conclude the presentation of our \textit{Morpheus} formulation by discussing the resulting trajectory for a head height reduction of $\Delta h = \SI{0.5}{\m}$.

\subsection{Resulting Motion} \label{sec:morpheus:results}
\begin{table}[h]
  \centering
  \begin{tabularx}{0.7\textwidth}{rcl}
  \toprule
    Symbol & Value & \multicolumn{1}{X}{Description} \\
  \midrule
    $\Delta h_\text{diff}$ & $\SI{0.5}{\metre}$ & Head Reduction \\
  \midrule
  PoS & \textit{\cref{eq:morpheus:constraints:stability}} & \acrlong{pos} \\
    $\gamma$ & $0.8$ & \acrshort{pos} safety Limit \\
    \textbf{ZMP} & \textit{\cref{eq:morpheus:constraints:stability}} & \acrlong{zmp} \\
  \midrule
    $\bm q$ & \textit{\cref{eq:morpheus:dyn:x}} & Joint Positions \\
  $\bm{\tau}$ & \textit{\cref{eq:morpheus:dyn:u}} & Joint Torques \\
    \midrule
  $\bm{\underline q},\ \bm{\overline q}$ & \textit{\cref{fig:morpheus:results:graphs}} & Joint Limits \\
    $\bm{\tau}_{\text{max}}$ & \textit{\cref{fig:morpheus:results:graphs}} & Max. Joint Torques \\
  \midrule
  $\bm F_j$ & \textit{\cref{eq:morpheus:constraints:F,eq:morpheus:constraints:Fz}} & Ground Reaction Forces \\
    $\mu$ & $0.5$ & Friction Coefficient \\
    $\alpha$ & $10^{-12}$ & Proximal Setting \\
  \midrule
    $v_{\text{stat.}}$ & \SI{0.1}{\metre \per \second} & Allowed Stationary Velocity \\
  \midrule
    $T_f$ & $\SI{1.5}{\second}$ & Terminal Time \\
    $N$ & $48$ & Shooting Nodes \\
    $m_h$ & \textit{\cref{eq:morpheus:obj:head}} & Mayer Obj. Head Term \\
    $\omega_h$ & $N$ & Mayer Obj. Head Term Weight \\
    $l_s$ & \textit{\cref{eq:morpheus:obj:stability}} & Lagrange Obj. Stability Term \\
    $\omega_s$ & $10^{-4}$ & Lagrange Obj. Stability Term Weight \\
    $l_r$ & \textit{\cref{eq:morpheus:obj:regularization}} & Lagrange Obj. Regularization Term \\
    $\omega_r$ & $10^{-8}$ & Lagrange Obj. Regularization Term Weight \\
  \bottomrule
  \end{tabularx}
  \caption{Variables and Parameters}
  \label{tab:morpheus:variables}
\end{table}
The collision avoidance motion in this section were generated for a duration of motion $T_f=\SI{1.5}{\second}$ with $N=48$ shooting nodes.
The head of our H1-2 has to be lowered by $\Delta h = \SI{0.5}{\meter}$, while an end velocity of $\bm v_\text{stat.}=\SI{0.1}{\radian \per \second}$ is allowed.
Additionally a \acrshort{pos} reduction of $\gamma=0.8$ is set to concede a safety margin.
To solve the forward constraint dynamics a proximal setting $\alpha=10^{-12}$ with ten unrolled iterations should lead to a close approximation of the original Jacobian.
A full list of all used parameters is displayed in \cref{tab:morpheus:variables}.
\\
In acados we utilize \acrshort{hpipm} \cite{frison_hpipm_2020}, as it is specifically build for abusing typical structures in robotics optimal control problems.
Underneath an explicit fourth order Runge-Kutta integrator is used.
To overcome strongly constraint regions, the line-search funnel-penalty method \cite{kiessling_2025} is the globalization strategy of choice.
It should be noted that classic merit backtracking did not converge for the given formulation.
\par
With this setup the \textit{Morpheus} algorithm achieved the motion displayed in \cref{fig:morpheus:results:banner}.
The figure was generated by capturing every tenth frame, starting from the first one and ending on the last available frame.
\begin{figure}
  \centering
  \makebox[\textwidth]{
  \centering
  \begin{subfigure}[b]{0.65\textwidth}
    \includegraphics[width=\textwidth]{figures/03_morpheusEvasion/results/joint6.pdf}
    \caption{Torso (Joint 6)}
    \label{fig:morpheus:results:graphs:torso}
  \end{subfigure}
  }
  \\
  \makebox[\textwidth]{
  \begin{subfigure}[b]{0.65\textwidth}
    \includegraphics[width=\textwidth]{figures/03_morpheusEvasion/results/joint0.pdf}
    \caption{Left Hip Pitch (Joint 0)}
    \label{fig:morpheus:results:graphs:leftHip}
  \end{subfigure}
  \hfill
  \begin{subfigure}[b]{0.65\textwidth}
    \includegraphics[width=\textwidth]{figures/03_morpheusEvasion/results/joint3.pdf}
    \caption{Right Hip Pitch (Joint 3)}
    \label{fig:morpheus:results:graphs:rightHip}
  \end{subfigure}
  }
  \\
  \makebox[\textwidth]{
  \begin{subfigure}[b]{0.65\textwidth}
    \includegraphics[width=\textwidth]{figures/03_morpheusEvasion/results/joint1.pdf}
    \caption{Left Knee (Joint 1)}
    \label{fig:morpheus:results:graphs:leftKnee}
  \end{subfigure}
  \hfill
  \begin{subfigure}[b]{0.65\textwidth}
    \includegraphics[width=\textwidth]{figures/03_morpheusEvasion/results/joint4.pdf}
    \caption{Right Knee (Joint 4)}
    \label{fig:morpheus:results:graphs:rightKnee}
  \end{subfigure}
  }
  \\
  \makebox[\textwidth]{
  \begin{subfigure}[b]{0.65\textwidth}
    \includegraphics[width=\textwidth]{figures/03_morpheusEvasion/results/joint2.pdf}
    \caption{Left Ankle Pitch (Joint 2)}
    \label{fig:morpheus:results:graphs:leftAnkle}
  \end{subfigure}
  \hfill
  \begin{subfigure}[b]{0.65\textwidth}
    \includegraphics[width=\textwidth]{figures/03_morpheusEvasion/results/joint5.pdf}
    \caption{Right Ankle Pitch (Joint 5)}
    \label{fig:morpheus:results:graphs:rightAnkle}
  \end{subfigure}
  }

  \caption{
    Joint trajectories for angles and torques, without floating base.
  }
  \label{fig:morpheus:results:graphs}
\end{figure}
All joint angles $\bm q_i$ as well as torques $\bm \tau_i$ are plotted with there hardware limits in \cref{fig:morpheus:results:graphs}.
\par
The reader will quickly notice, that, as a consequence of our contact frame definition, the left and right side trajectories for joint angles as well as torques are nearly identical.
Small divergence in torque is probably due to the piece-wise constant torque control.
\\
While neither the available range of motion nor the maximum torques were a limitation for the torso joint (\cref{fig:morpheus:results:graphs:torso}), a clear indication towards such constraints can be found for the ankles and hips (\cref{fig:morpheus:results:graphs:leftAnkle,fig:morpheus:results:graphs:rightAnkle,fig:morpheus:results:graphs:leftHip,fig:morpheus:results:graphs:rightHip}).
As it is the case for humans, the natural range of motion for the H1-2, is most limited in pitch direction at the ankle, compared to knee and hip joints.
If we compare the torque limits of ankle and knee joints multi-joint torques for humans, a shortcoming in the ankle area can be found.
Hahn et al. \cite{hahn_2011} recorded a $\SI{289.57 \pm 43.3}{\N \m}$ at \SI{0.88 \pm 0.16}{\radian} for the knee joint and $\SI{234.57 \pm 44.0}{\N \m}$ for ankle pitch joint during a multi-joint action.
Unitrees H1-2 seems to match the peak torque capabilities of the human body in the knee, while falling short of approximately $\SI{100}{\N \m}$ at the ankle joint.
However the presented peak torque values for the human ankle joint can only be reached in a range between $\SI{-0.14 \pm 0.04}{\radian}$ and $\SI{0.09 \pm 0.04}{\radian}$.
Outside this window, Hahn et al. report a torque limit of $\SI{130.3 \pm 25.5}{\N \m}$ at $\SI{0.40 \pm 0.04}{\radian}$.
This lowered peak torque matches the capabilities of the H1-2, for single pitch joint usage.
It is therefore apparent, that while Unitrees H1-2 has stronger knee joints, its ankles are a weak point, especially, when it comes to motions where the feet are mostly orthogonal to the lower shank.
The result is, that some motions, which are easy to execute for humans pose problems for the H1-2 hardware.
In our case however is the motion mostly limited in areas where a human would also struggle to develop enough strength.
At then end of the motion the robot is bend over backwards, generating requiring high torques, with at its joint angle limits.
\\
This shows that the combination of our \textit{Morpheus} generator with the H1-2 Unitree humanoid is well suited to find human like motions.
\par
An indicator of the hardware limitation is the generated motion and its joint trajectories itself.
\Cref{fig:morpheus:results:banner} shows, that the generated motion abuses the fact, that we do not set head height requirements for the time until it reaches its final pose.
The result is a swing up move in the middle of the motion, which can be observed in the approximate time range from $\SI{0.25}{\s}$ to $\SI{1.0}{\s}$.
Because of the regularization term $l_r$ (\cref{eq:morpheus:obj:regularization}) and torque limits (\cref{eq:morpheus:constraints:joint}) its is probably more feasible to find a solution that keeps the robot in motion and preserve momentum instead of holding constant torque on the motors for a leaned over pose.
Therefore the peak in ankle torques primarily occurs at the end of the motion, when the swing has to be stopped for the final quasi-static pose.
\\
Our motion synthesis formulation is able to successfully produce suitable whole-body motions to push its head underneath a specified height, while exploring the capabilities of the provided hardware.
However, the swing motion shows the need for a dodge constraints, which not only imposes constraints on the last shooting node, but enforces a safety distance between approaching object and the robots body.
